\documentclass[11pt,a4paper,oneside]{report}             % Single-side
%\documentclass[11pt,a4paper,twoside,openright]{report}  % Duplex

%\PassOptionsToPackage{chapternumber=Huordinal}{magyar.ldf}
\usepackage{t1enc}
\usepackage[utf8]{inputenc}
\usepackage{amsmath}
\usepackage{amssymb}
\usepackage{enumerate}
\usepackage[thmmarks]{ntheorem}
\usepackage{graphics}
\usepackage{epsfig}
\usepackage{listings}
\usepackage{color}
%\usepackage{fancyhdr}
\usepackage{lastpage}
\usepackage{anysize}
\usepackage[magyar,english]{babel}
\usepackage{sectsty}
\usepackage{setspace}  % Ettol a tablazatok, abrak, labjegyzetek maradnak 1-es sorkozzel!
\usepackage[hang]{caption}
\usepackage{hyperref}

%--------------------------------------------------------------------------------------
% Main variables
%--------------------------------------------------------------------------------------
\newcommand{\vikszerzo}{Kovács Ádám}
\newcommand{\vikszerzos}{Gémes Andrea Kinga}
\newcommand{\vikkonzulens}{Dr.~Recski Gábor}
\newcommand{\vikcim}{Semantic parsing with graph transformations}
\newcommand{\viktanszek}{Department of Automation and Applied Informatics}
\newcommand{\vikdoktipus}{Scienfitic Students' Assosiactions report}
\newcommand{\vikdepartmentr}{Kovács Ádám}

%--------------------------------------------------------------------------------------
% Page layout setup
%--------------------------------------------------------------------------------------
% we need to redefine the pagestyle plain
% another possibility is to use the body of this command without \fancypagestyle
% and use \pagestyle{fancy} but in that case the special pages
% (like the ToC, the References, and the Chapter pages)remain in plane style

\pagestyle{plain}
%\setlength{\parindent}{0pt} % ?ttekinthet?bb, angol nyelv? dokumentumokban jellemz?
%\setlength{\parskip}{8pt plus 3pt minus 3pt} % ?ttekinthet?bb, angol nyelv? dokumentumokban jellemz?
\setlength{\parindent}{12pt} % magyar nyelv? dokumentumokban jellemz?
\setlength{\parskip}{0pt}    % magyar nyelv? dokumentumokban jellemz?

\marginsize{35mm}{25mm}{15mm}{15mm} % anysize package
\setcounter{secnumdepth}{0}
\sectionfont{\large\upshape\bfseries}
\setcounter{secnumdepth}{2}
\singlespacing
\frenchspacing

%--------------------------------------------------------------------------------------
%	Setup hyperref package
%--------------------------------------------------------------------------------------
\hypersetup{
    bookmarks=true,            % show bookmarks bar?
    unicode=false,             % non-Latin characters in Acrobat?s bookmarks
    pdftitle={\vikcim},        % title
    pdfauthor={\vikszerzo},    % author
    pdfsubject={\vikdoktipus}, % subject of the document
    pdfcreator={\vikszerzo},   % creator of the document
    pdfproducer={Producer},    % producer of the document
    pdfkeywords={keywords},    % list of keywords
    pdfnewwindow=true,         % links in new window
    colorlinks=true,           % false: boxed links; true: colored links
    linkcolor=black,           % color of internal links
    citecolor=black,           % color of links to bibliography
    filecolor=black,           % color of file links
    urlcolor=black             % color of external links
}

%--------------------------------------------------------------------------------------
% Set up listings
%--------------------------------------------------------------------------------------
\lstset{
	basicstyle=\scriptsize\ttfamily, % print whole listing small
	keywordstyle=\color{black}\bfseries\underbar, % underlined bold black keywords
	identifierstyle=, 					% nothing happens
	commentstyle=\color{white}, % white comments
	stringstyle=\scriptsize\sffamily, 			% typewriter type for strings
	showstringspaces=false,     % no special string spaces
	aboveskip=3pt,
	belowskip=3pt,
	columns=fixed,
	backgroundcolor=\color{lightgray},
} 		
\def\lstlistingname{lista}	

%--------------------------------------------------------------------------------------
%	Some new commands and declarations
%--------------------------------------------------------------------------------------
\newcommand{\code}[1]{{\upshape\ttfamily\scriptsize\indent #1}}

% define references
\newcommand{\figref}[1]{\ref{fig:#1}.}
\renewcommand{\eqref}[1]{(\ref{eq:#1})}
\newcommand{\listref}[1]{\ref{listing:#1}.}
\newcommand{\sectref}[1]{\ref{sect:#1}}
\newcommand{\tabref}[1]{\ref{tab:#1}.}

\DeclareMathOperator*{\argmax}{arg\,max}
%\DeclareMathOperator*[1]{\floor}{arg\,max}
\DeclareMathOperator{\sign}{sgn}
\DeclareMathOperator{\rot}{rot}
\definecolor{lightgray}{rgb}{0.95,0.95,0.95}

\author{\vikszerzo \and \vikszerzos}
\title{\viktitle}
\includeonly{
	titlepage,%
	declaration,%
	abstract,%
	introduction,%
	chapter1,%
	chapter2,%
	chapter3,%
	acknowledgement,%
}
%--------------------------------------------------------------------------------------
%	Setup captions
%--------------------------------------------------------------------------------------
\captionsetup[figure]{
%labelsep=none,
%font={footnotesize,it},
%justification=justified,
width=.75\textwidth,
aboveskip=10pt}

\renewcommand{\captionlabelfont}{\small\bf}
\renewcommand{\captionfont}{\footnotesize\it}

%--------------------------------------------------------------------------------------
% Table of contents and the main text
%--------------------------------------------------------------------------------------
\begin{document}
\singlespacing
\include{guideline}
\include{project}

\pagenumbering{arabic}
\onehalfspacing
%--------------------------------------------------------------------------------------
%	The title page
%--------------------------------------------------------------------------------------
\begin{titlepage}
\begin{center}
\includegraphics[width=60mm,keepaspectratio]{figures/BMElogo.png}\\
\vspace{0.3cm}
\textbf{Budapesti Műszaki és Gazdaságtudományi Egyetem}\\
\textmd{Villamosmérnöki és Informatikai Kar}\\
\textmd{\viktanszek}\\[5cm]

\vspace{0.4cm}
{\huge \bfseries \vikcim}\\[0.8cm]
\vspace{0.5cm}
\textsc{\Large \vikdoktipus}\\[4cm]

\begin{tabular}{cc}
 \makebox[7cm]{\emph{Author}} & \makebox[7cm]{\emph{Supervisor}} \\
 \makebox[7cm]{\vikszerzo} & \makebox[7cm]{\vikkonzulens} \\
 \makebox[7cm]{\vikszerzos}
\end{tabular}

\vfill
{\large \today}
\end{center}
\end{titlepage}



\tableofcontents\vfill
%----------------------------------------------------------------------------
% Abstract in hungarian
%----------------------------------------------------------------------------
\chapter*{Kivonat}\addcontentsline{toc}{chapter}{Kivonat}
%--------------------------------------------------------------------------------------
% semantic parsing, graftranszformacio altalaban
A szemantikai elemzés célja, hogy természetes nyelvi adathoz
készíthessünk szemantikai reprezentációt, így tudjuk modellezni a szöveg jelentését. Ha a nyelvi
jelentést fogalmak irányított gráfjaival reprezentáljuk, ezeket pedig a mondat szintaktikai
szerkezetét reprezentáló fákból kell előállítanunk, akkor a teljes feladat egyetlen komplex gráftranszformációként
definiálható.

%--------------------------------------------------------------------------------------
% machine comprehension
A természetes nyelv szemantikájának reprezentációját ritkán használják közvetlenül a state-of-the art rendszerekben népszerű szemantikai feladatokra, mint a szemantikai hasonlóság mérése, vagy a gépi szövegértés. Ezek a rendszerek többnyire szó embeddingeket használnak a szavak jelentésének ábrázolására.
Ebben a dolgozatban mi gráf reprezentációkat és transzformációkat használunk, mint egyszerű ám hatékony eszközök az oksági viszony felismerésére, valamint leírunk egy módszert a \texttt{4lang} szemantikus elemzőrendszer (Recski:2016) használatára a 2018-as Semeval Task \textit{Machine comprehension using commonsense knowledge} kapcsán. Ez a feladat azt kívánja a résztvevőktől, hogy olyan rendszereket tanítsanak fel, amelyek ki tudják választani a megfelelő választ az egyszerű, több válaszlehetőséget kínáló kérdéseknél rövid eseményleíró szövegek elolvasása után. A tanító és teszt adat az \texttt{MCScript} adathalmaz (Ostermann et al., 2018) részhalmazából lett kinyerve.
A két legjobb rendszer, \texttt{HFL-RC} (Chen et al., 2018) és \texttt{Yuanfudao} egyenként $84,15\%$ és $83,95\%$ pontosságot ért el a teszt adaton.

%--------------------------------------------------------------------------------------
% baseline
Először bemutatunk egy hatékony baselinet ezen a feladaton csupán a szemantikus gráfok és a köztük lévő hasonlóságok felhasználásával. Ezt követően leírjuk a \texttt{Yuanfudao} (Wang et al., 2018) state-of-the art rendszert és az ezzel végzett kísérletezéseinket, amelyek során
%--------------------------------------------------------------------------------------
% rendszer, amivel foglalkoztunk
a baselineunkat extra featureként felhasználva javítottunk a neurális hálón. A rendszer kiválasztása magától értetődő volt, mivel a forráskód nyilvánosan elérhető és már sikeresen alkalmazott tudás alapú reprezentációt a szópárok közötti szemantikai kapcsolatokra, a \textit{ConceptNet}et.
Eredményeink azt mutatják, hogy ezzel a módosítással $0,5\%$ százalékpont növekedés érhető el és a \textit{ConceptNet} helyettesíthető a mi szemantikus modellünkkel.
\vfill

%----------------------------------------------------------------------------
% Abstract in english
%----------------------------------------------------------------------------
\chapter*{Abstract}\addcontentsline{toc}{chapter}{Abstract}
%--------------------------------------------------------------------------------------
% semantic parsing, graftranszformacio altalaban
The main task of semantic parsing is to automatically build semantic representation from the input, so we can model the meaning
of raw texts. If we model meaning as directed graphs of concept and we can build them from syntax trees that represent the structure of 
sentences, then we can define the whole process as one complex graph transformation.

%--------------------------------------------------------------------------------------
% machine comprehension
Representations of natural language semantics are rarely used explicitly in state-of-the art systems for popular semantics tasks such as measuring semantic similarity or machine comprehension. These systems mostly use word embeddings as representation of word meaning.

In this paper we use graphical representations and transformations as simple but powerful tools for recognizing entailment and we describe a method using semantic parsing system \texttt{4lang} (Recski:2016) and applying it on the 2018 Semeval Task \textit{Machine comprehension using commonsense knowledge}. This task requires participants to train systems that can choose the correct answer to simple multiple choice questions based on short passages describing simple chains of events. Data for both training and testing is extracted from the \texttt{MCScript} dataset (Ostermann et al., 2018). The top two systems, \texttt{HFL-RC} (Chen et al., 2018) and \texttt{Yuanfudao} achieved accuracy scores of $84.15\%$ and $83.95\%$ on the test data, respectively.

%--------------------------------------------------------------------------------------
% baseline
First we will present a strong baseline on this task using only semantic graphs and similarities among them, followed by 
a description of a state-of-the art system \texttt{Yuanfudao} (Wang et al., 2018) and our experiments with it where
%--------------------------------------------------------------------------------------
% rendszer, amivel foglalkoztunk
we used our baseline as an extra feature for improving the neural network. The choice of the system was obvious, because the source code is publicly available and it already employs successfully a knowledge base representing semantic relationships among pairs of words, \textit{ConceptNet}.
Our results suggest that these features achieve a .5 percentage point improvement, and the \textit{ConceptNet} could be replaced by our semantic model.
\vfill


\chapter{Introduction}
\label{chap:Introdu}
%----------------------------------------------------------------------------
\section{Natural Language Processing}
While computers can be easily programmed to understand structured data, such as tables and spreadsheets, it can be rather challenging for them
to understand human communication. Because there is a vast difference in the magnitude of the unstructured data compared to the structured ones, there is a high demand
for tools, that can deal with raw text. Thats where NLP comes in. It contains high variety of tools, that we have to use, when we need to deal with natural input.

Every day, we come into contact with human communication, we say a lot of words to other people, and they try to interpret them even when the context of the saying 
isn't necesseraily complete. The listeners can use their common knowledge to fill the needed information. We resolve ambiguities, misunderstanding, and can even understand words 
we have never heard before just from the context of the communication.
Even though these tasks are trivial for us, for a computer it can be really hard.

%--------------------------------------------------------------------------------------
% egyszer� multtol jelenig bevezeto
The interest in NLP research began in the 1950s, the early phase was mainly focused on MT (Machine Translation), because after the World War II, people
recognised the importance of the translation from one language to another, and hoped to do it automatically.
However MT is still a very difficult nowadays, so these researches discovered early the main challenges of the syntactic and semantic parsing.
As time passed, researches embraced new areas of NLP as more advanced technology and knowledge became available. Now that we live in a world where
computers and smartphones, collecting data became incredibly easy, as a result, statistical NLP drew attention because these models thrive off big data, but one cannot ignore 
simple rule based methods, which can also be a very powerful method, especially using them as a hybrid model with statistical methods.
%--------------------------------------------------------------------------------------
% r�vid nlp pipeline

Building NLP applications requires many levels of analysis.
The typical pipeline structured as follows:
\begin{itemize}
	\item First we need to tokenize our input text, which means breaking up the text into meaningful elements, especially into words
	\item After we tokenized our text, we need word analysis called \texttt{morphology}, which is concerned with the structure of words.
	\item Part of speech explains how a word is used in a sentence, and tagging is a process marking words to its particular part of speech.
	\item The main task of syntactic parsing is to analyze the grammatical structure of a sentence. Given a set of words, a parser forms units (subjects, verbs, etc..) according to some grammatical formalism.
	There are two main type of syntactic parsers:
	\begin{itemize}
		\item Consituency parsers produce trees, that represents the grammatical structure.
		\item Dependency parsers are the more popular nowadays. They represent the structure of a sentence as a dependency tree.
	\end{itemize}
	\item At this point we have various ways to analyze a text, but without modeling its \texttt{meaning}. Semantic is the study of meaning, and semantic parsing is a task to find a representation and assign it to the text. This task will be the main topic of my work.
\end{itemize}

%\listoffigures\addcontentsline{toc}{chapter}{?br?k jegyz?ke}
%\listoftables\addcontentsline{toc}{chapter}{T?bl?zatok jegyz?ke}

%\bibliography{mybib}
%\addcontentsline{toc}{chapter}{Irodalomjegyz?k}
\bibliography{mybib}
\addcontentsline{toc}{chapter}{References}
\bibliographystyle{plain}

\include{appendices}

\label{page:last}
\end{document}
